%! Equivalent Representations of Trigonometric Functions — Unit 7, Lesson 7.5 — Group Activity (Answer Key)
\documentclass[10pt]{article}
\usepackage{precalculus-article}
\usepackage{precalculus-key}   % pulls in -boxes; do NOT also load -boxes
\newcommand{\TallMath}[1]{$\displaystyle #1\rule[-1.4em]{0pt}{3.2em}$}

\begin{document}

\pageheader{Unit 7, Lesson 7.5}{Group Activity --- Answer Key}
\namepartnerperiod

\vspace{0.10in}

\begin{scenariobox}[Context: Solving Equations Using Trigonometric Identities]{plumlight}
Equations like $2\sin^2 x + \cos x = 1$ mix two different trig functions and
cannot be solved directly. In this activity your group will use the Pythagorean,
double-angle, and sum identities to rewrite such equations in terms of a
\emph{single} trig function, then solve them algebraically.
\end{scenariobox}

\vspace{0.08in}

%----------------------------------------------------------------------
% Tier R — Key
%----------------------------------------------------------------------
\begin{tcolorbox}[colback=white, colframe=black!40,
    title=\textbf{Tier R --- Reinforce (Key)},
    fonttitle=\bfseries, arc=2mm, left=3mm, right=3mm, top=2mm, bottom=2mm]

\textbf{Step 1.} Substitute $\sin^2 x = 1 - \cos^2 x$:

$2(\;$ \ans{$1 - \cos^2 x$} $\;) + \cos x - 1 = 0$

\vspace{6pt}
\textbf{Step 2.} Expand: \ans{$2 - 2\cos^2 x + \cos x - 1 = 0$}

\vspace{6pt}
\textbf{Step 3.} Collect and multiply by $-1$: \ans{$2\cos^2 x - \cos x - 1 = 0$}

\vspace{6pt}
\textbf{Step 4.} Let $u = \cos x$: \ans{$2u^2 - u - 1 = 0$}

Factor: \ans{$(2u + 1)(u - 1) = 0$}

\vspace{6pt}
\textbf{Step 5.} $\cos x = $ \ans{$-\dfrac{1}{2}$} \qquad or \qquad $\cos x = $ \ans{1}

\vspace{6pt}
\textbf{Step 6.}

\renewcommand{\arraystretch}{1.8}
\begin{tabular}{|l|l|}
\hline
$\cos x = $ \ans{$-\tfrac{1}{2}$} & Solutions: \ans{$\dfrac{2\pi}{3},\;\dfrac{4\pi}{3}$} \\
\hline
$\cos x = $ \ans{1} & Solutions: \ans{0} \\
\hline
\end{tabular}

\vspace{6pt}
\textbf{Answer:} $x \in \big\{$ \ans{$0,\;\dfrac{2\pi}{3},\;\dfrac{4\pi}{3}$} $\big\}$

\vspace{6pt}
\textbf{Check:} $x = 0$: $2\sin^2(0) + \cos(0) - 1 = $ \ans{0} $+$ \ans{1} $- 1 =$ \ans{0} \checkmark

\end{tcolorbox}

\vspace{0.08in}

%----------------------------------------------------------------------
% Tier A — Key
%----------------------------------------------------------------------
\begin{tcolorbox}[colback=white, colframe=black!40,
    title=\textbf{Tier A --- Apply (Key)},
    fonttitle=\bfseries, arc=2mm, left=3mm, right=3mm, top=2mm, bottom=2mm]

\textbf{Problem 1.} $2\cos^2 x - \cos x - 1 = 0$

Identity used: \ans{None needed (already one function)}

Let $u = \cos x$: $(2u + 1)(u - 1) = 0 \Rightarrow \cos x = -\tfrac{1}{2}$ or $\cos x = 1$.

\textbf{Solutions:} $x \in \big\{$ \ans{$0,\;\dfrac{2\pi}{3},\;\dfrac{4\pi}{3}$} $\big\}$

\vspace{8pt}
\textbf{Problem 2.} $\sin 2x = \sin x$

Identity used: \ans{Double-angle: $\sin 2x = 2\sin x\cos x$}

$2\sin x\cos x = \sin x \Rightarrow 2\sin x\cos x - \sin x = 0 \Rightarrow \sin x(2\cos x - 1) = 0$

$\sin x = 0$: $x = 0, \pi$ \quad or \quad $\cos x = \tfrac{1}{2}$: $x = \tfrac{\pi}{3}, \tfrac{5\pi}{3}$

\textbf{Solutions:} $x \in \big\{$ \ans{$0,\;\dfrac{\pi}{3},\;\pi,\;\dfrac{5\pi}{3}$} $\big\}$

\vspace{8pt}
\textbf{Problem 3.} $\cos^2 x - \sin^2 x = \dfrac{1}{2}$

Identity used: \ans{Double-angle: $\cos 2x = \cos^2 x - \sin^2 x$}

$\cos 2x = \tfrac{1}{2} \Rightarrow 2x = \tfrac{\pi}{3}, \tfrac{5\pi}{3}, \tfrac{7\pi}{3}, \tfrac{11\pi}{3}$

$\Rightarrow x = \tfrac{\pi}{6}, \tfrac{5\pi}{6}, \tfrac{7\pi}{6}, \tfrac{11\pi}{6}$

\textbf{Solutions:} $x \in \big\{$ \ans{$\dfrac{\pi}{6},\;\dfrac{5\pi}{6},\;\dfrac{7\pi}{6},\;\dfrac{11\pi}{6}$} $\big\}$

\end{tcolorbox}

\vspace{0.08in}

%----------------------------------------------------------------------
% Tier E — Key
%----------------------------------------------------------------------
\begin{tcolorbox}[colback=white, colframe=black!40,
    title=\textbf{Tier E --- Extend (Key)},
    fonttitle=\bfseries, arc=2mm, left=3mm, right=3mm, top=2mm, bottom=2mm]

\textbf{Part A --- Verify Identities.}

\textbf{Identity 1:} $\dfrac{\sin^2\theta}{1 - \cos\theta} = 1 + \cos\theta$

\ans{Left side: $\dfrac{1 - \cos^2\theta}{1 - \cos\theta} = \dfrac{(1-\cos\theta)(1+\cos\theta)}{1-\cos\theta} = 1 + \cos\theta$ = Right side. $\checkmark$}

\vspace{6pt}
\textbf{Identity 2:} $\tan\theta + \cot\theta = \dfrac{1}{\sin\theta\cos\theta}$

\ans{Left side: $\dfrac{\sin\theta}{\cos\theta} + \dfrac{\cos\theta}{\sin\theta} = \dfrac{\sin^2\theta + \cos^2\theta}{\sin\theta\cos\theta} = \dfrac{1}{\sin\theta\cos\theta}$ = Right side. $\checkmark$}

\vspace{8pt}
\textbf{Part B --- Exact Value.}

$\sin\!\left(\dfrac{5\pi}{12}\right) = \sin\!\left(\dfrac{\pi}{4} + \dfrac{\pi}{6}\right)
= \sin\dfrac{\pi}{4}\cos\dfrac{\pi}{6} + \cos\dfrac{\pi}{4}\sin\dfrac{\pi}{6}$

$= \dfrac{\sqrt{2}}{2}\cdot\dfrac{\sqrt{3}}{2} + \dfrac{\sqrt{2}}{2}\cdot\dfrac{1}{2}
= \dfrac{\sqrt{6}}{4} + \dfrac{\sqrt{2}}{4}$

$\sin\!\left(\dfrac{5\pi}{12}\right) = $ \ans{$\dfrac{\sqrt{6}+\sqrt{2}}{4}$}

\end{tcolorbox}

\end{document}
